\documentclass{beamer}
\usetheme{Madrid}

% Encoding and fonts for pdfLaTeX
\usepackage[utf8]{inputenc}
\usepackage[T1]{fontenc}
\usepackage{lmodern}
\usepackage{hyperref}

\title{Prompt Engineering}
\author{Mehmet Ali Akyol, PhD}
\institute{Ankara Medipol University}
\date{\today}

\begin{document}

\begin{frame}
    \titlepage
    \begin{center}
        \vspace{1cm}
        \textbf{Week 8: Evaluating AI Outputs}
    \end{center}
\end{frame}

\section{Overview}

\begin{frame}
    \frametitle{The Evaluation Gap}
    \begin{block}{The Paradox of Generative AI}
        It is significantly easier to \textbf{generate} text than to \textbf{verify} it.
    \end{block}
    \begin{itemize}
        \item \textbf{Generation:} Low effort, high speed, infinite variability.
        \item \textbf{Verification:} High effort, requires domain expertise, slow.
        \item \textbf{The Risk:} Without rigorous evaluation, we scale errors, not solutions.
        \item \textbf{The Goal:} Close the gap by automating verification where possible and focusing human effort where necessary.
    \end{itemize}
\end{frame}

\begin{frame}
    \frametitle{Why Evaluation Matters}
    \begin{itemize}
        \item \textbf{Reliability:} Moving from ``it looks good'' to ``it works 99\% of the time.''
        \item \textbf{Safety:} Detecting edge cases, jailbreaks, and harmful outputs before deployment.
        \item \textbf{Iterative Improvement:} You cannot improve what you cannot measure. Evaluation provides the signal for the refinement loop.
        \item \textbf{Model Selection:} Data-driven decisions on which model (GPT-4 vs. Llama 3) offers the best cost/performance ratio.
    \end{itemize}
\end{frame}

\begin{frame}
    \frametitle{Weekly Learning Objectives}
    \begin{itemize}
        \item Define evaluation criteria aligned with task goals.
        \item Apply qualitative and quantitative assessment methods.
        \item Design human-in-the-loop review pipelines.
        \item Use automation (judge models, checklists) responsibly.
        \item Interpret evaluation results to guide refinement.
    \end{itemize}
\end{frame}

\begin{frame}
    \frametitle{Key Concepts}
    \begin{itemize}
        \item Quality dimensions (accuracy, completeness, tone, compliance)
        \item Intrinsic vs. extrinsic evaluation
        \item Rubrics, scorecards, and annotation protocols
        \item Sampling strategies (random, risk-based, stratified)
        \item Inter-rater reliability and bias mitigation
    \end{itemize}
\end{frame}

\section{Designing Evaluation Criteria}

\begin{frame}
    \frametitle{Framework: The 3 H's of Evaluation}
    \begin{columns}[T]
        \column{0.33\textwidth}
        \textbf{1. Helpful}
        \begin{itemize}
            \item Does it solve the user's problem?
            \item Is the format correct?
            \item Is the tone appropriate?
        \end{itemize}
        
        \column{0.33\textwidth}
        \textbf{2. Honest}
        \begin{itemize}
            \item Is it factually accurate?
            \item Does it admit uncertainty?
            \item No hallucinations?
        \end{itemize}
        
        \column{0.33\textwidth}
        \textbf{3. Harmless}
        \begin{itemize}
            \item No toxicity or bias.
            \item No dangerous content.
            \item Resists jailbreaks.
        \end{itemize}
    \end{columns}
\end{frame}

\begin{frame}
    \frametitle{From Task to Metrics}
    \begin{itemize}
        \item \textbf{Clarify Intent:} What is the ``Job to be Done''? (e.g., Summarize vs. Analyze).
        \item \textbf{Define Failure Modes:} What does a \textit{bad} response look like? (e.g., Missing citations, wrong JSON format).
        \item \textbf{Balance Metrics:}
        \begin{itemize}
            \item \textit{Hard Metrics:} Syntax, length, forbidden words (Deterministic).
            \item \textit{Soft Metrics:} Creativity, empathy, flow (Probabilistic).
        \end{itemize}
    \end{itemize}
\end{frame}

\begin{frame}
    \frametitle{Rubric Construction: Best Practices}
    \begin{itemize}
        \item \textbf{Binary Checklists:} Best for strict constraints.
        \begin{itemize}
            \item ``Does the output contain a disclaimer?'' (Yes/No)
        \end{itemize}
        \item \textbf{Likert Scales (1--5):} Best for nuance.
        \begin{itemize}
            \item 1 = Completely wrong.
            \item 3 = Mostly correct but requires editing.
            \item 5 = Perfect, ready to publish.
        \end{itemize}
        \item \textbf{Comparative Ranking:} Best for preference.
        \begin{itemize}
            \item ``Is Output A better than Output B?''
        \end{itemize}
    \end{itemize}
\end{frame}

\section{Evaluation Methods}

\begin{frame}
    \frametitle{The Golden Dataset (Ground Truth)}
    \begin{block}{Definition}
        A curated set of inputs and their corresponding ``perfect'' outputs, used as a benchmark for testing.
    \end{block}
    \begin{itemize}
        \item \textbf{Creation:} Hand-written by experts or verified high-quality model outputs.
        \item \textbf{Size:} Start small (10--50 examples) but high diversity.
        \item \textbf{Usage:} Run every prompt version against this set to detect regression.
        \item \textbf{Maintenance:} Update the dataset as new edge cases are discovered in production.
    \end{itemize}
\end{frame}

\begin{frame}
    \frametitle{Qualitative Techniques (Human-Centric)}
    \begin{itemize}
        \item \textbf{Expert Review:} Subject Matter Experts (SMEs) grade outputs. High cost, high accuracy.
        \item \textbf{Crowdsourcing:} Non-experts rate for general feel/fluency. Lower cost, variable quality.
        \item \textbf{Red Teaming:} Adversarial testing to find safety flaws.
        \item \textbf{User Feedback:} Thumbs up/down, ``Regenerate'' rate in production.
    \end{itemize}
\end{frame}

\begin{frame}
    \frametitle{Quantitative Techniques (Metric-Centric)}
    \begin{itemize}
        \item \textbf{N-gram Metrics (BLEU, ROUGE):} Measures word overlap. Good for translation, bad for reasoning.
        \item \textbf{Embedding Similarity:} Measures semantic closeness to the Golden Dataset using vector distance (Cosine Similarity).
        \item \textbf{Code Execution:} Does the generated code run? Does it pass unit tests?
        \item \textbf{Factuality Scores:} Using search engines to verify claims (e.g., RAGAS metrics).
    \end{itemize}
\end{frame}

\begin{frame}
    \frametitle{Sampling Strategies}
    \begin{itemize}
        \item Random sampling for baseline quality tracking.
        \item Risk-based sampling (e.g., high-stakes domains, new templates).
        \item Stratified sampling across personas, use cases, data segments.
        \item Event-driven sampling after major prompt or model changes.
    \end{itemize}
\end{frame}

\section{Automation and Tooling}

\begin{frame}
    \frametitle{LLM-as-a-Judge Architectures}
    \begin{itemize}
        \item \textbf{Direct Scoring:} ``Rate this response from 1 to 5 based on clarity.''
        \item \textbf{Pairwise Comparison:} ``Which response is better, A or B?'' (Often more reliable than scoring).
        \item \textbf{Reference-Based:} ``Compare the response to this Golden Answer. Is it factually consistent?''
        \item \textbf{Rubric-Based:} ``Does the response meet criteria X, Y, and Z? Answer Yes/No for each.''
    \end{itemize}
\end{frame}

\begin{frame}
    \frametitle{Judge Models and Checklists}
    \begin{itemize}
        \item \textbf{The Judge:} A stronger model (e.g., GPT-4) evaluates a weaker model (e.g., Llama-7B).
        \item \textbf{Deterministic Checks:} Regex for phone numbers, JSON validation, forbidden word lists.
        \item \textbf{Hybrid Pipelines:}
        \begin{enumerate}
            \item Fast automated checks (Syntax).
            \item LLM Judge (Semantics).
            \item Human Review (Edge cases/Low confidence scores).
        \end{enumerate}
    \end{itemize}
\end{frame}

\begin{frame}
    \frametitle{Dashboards and Reporting}
    \begin{itemize}
        \item Track key metrics over time (accuracy, adherence, response time).
        \item Segment by task, persona, channel for targeted insights.
        \item Highlight trends, anomalies, and top failure modes.
        \item Share reports with stakeholders to drive accountability.
    \end{itemize}
\end{frame}

\section{Bias and Reliability}

\begin{frame}
    \frametitle{Common Evaluator Biases}
    \begin{itemize}
        \item \textbf{Verbosity Bias:} Models (and humans) tend to rate longer answers as better, even if they are fluff.
        \item \textbf{Egocentric Bias:} LLMs prefer outputs generated by themselves or models from the same family.
        \item \textbf{Position Bias:} In pairwise comparison, the first option (Option A) is often preferred.
        \item \textbf{Fix:} Swap positions (A/B testing), force brevity penalties, and use blind grading.
    \end{itemize}
\end{frame}

\begin{frame}
    \frametitle{Reducing Reviewer Bias}
    \begin{itemize}
        \item \textbf{Blind Reviews:} Hide the model name/version from the human reviewer.
        \item \textbf{Calibration Sessions:} Have all reviewers grade the same 5 examples and discuss discrepancies.
        \item \textbf{Rotation:} Rotate reviewers to prevent fatigue and ``anchoring'' to a specific style.
        \item \textbf{Gold Standard Injection:} Slip known ``bad'' examples into the queue to check if reviewers are paying attention.
    \end{itemize}
\end{frame}

\begin{frame}
    \frametitle{Inter-Rater Reliability}
    \begin{itemize}
        \item Measure agreement (Cohen's kappa, Krippendorff's alpha).
        \item Investigate disagreements to refine rubrics.
        \item Provide refresher training and updated exemplars.
        \item Automate alerts when reliability drops below threshold.
    \end{itemize}
\end{frame}

\section{Worked Examples}

\begin{frame}
    \frametitle{Legal Draft Review}
    \begin{itemize}
        \item Criteria: accuracy, completeness, tone, compliance.
        \item Process: human review + policy checklist + judge model.
        \item Report: monthly dashboard highlighting clause omissions.
        \item Outcome: 30\% reduction in manual escalations.
    \end{itemize}
\end{frame}

\begin{frame}
    \frametitle{Marketing Campaign QA}
    \begin{itemize}
        \item Criteria: brand voice, CTA clarity, localization quality.
        \item Use A/B testing metrics (CTR, conversion) for extrinsic evaluation.
        \item Sample high-risk markets at higher frequency.
        \item Share scorecards with creative team for continuous improvement.
    \end{itemize}
\end{frame}

\begin{frame}
    \frametitle{Academic Summaries}
    \begin{itemize}
        \item Criteria: factual accuracy, citation fidelity, coverage of sections.
        \item Workflow: human subject matter expert review + fact-check prompts.
        \item Use question-answer pairs to probe understanding.
        \item Publish evaluation logs with each delivered summary.
    \end{itemize}
\end{frame}

\section{Wrap-up}

\begin{frame}
    \frametitle{Key Takeaways}
    \begin{itemize}
        \item Evaluation turns AI outputs into dependable deliverables.
        \item Combine qualitative judgment with quantitative tracking.
        \item Automate selectively; keep humans where nuance matters.
        \item Use evaluation insights to steer prompt and process refinements.
    \end{itemize}
\end{frame}

\begin{frame}{Questions and Discussion}
    \begin{center}
    {\Huge Thank You!}
    \vspace{1cm}

        \textbf{Dr.\ Mehmet Ali Akyol}\\
        \texttt{mehmetali.akyol@gmail.com}
    \vspace{1cm}

    {\large Questions?}
    \end{center}
\end{frame}

\end{document}
